\documentclass[10pt]{article}
\usepackage{dcn}
\tikzset{
  pc/.style={hostn,minimum width=5mm,minimum height=4mm,inner sep=1pt,font=\tiny},
  term/.style={fill=ink,minimum width=1.2mm,minimum height=4mm,inner sep=0pt},
  tcap/.style={font=\small\bfseries,color=brand},
}
\newcommand{\figcap}[1]{{\scriptsize\color{midgrey}#1}}
\begin{document}

\packtitle{01}{Fundamentals: Components, Modes, Network Types \& Topologies}
  {Course outline weeks 1--2: concepts, goals \& applications, network structure, topology design}{HIGH}

This topic appears in \hl{every} paper, usually as 4--12 mark ``explain / differentiate / draw'' parts.
These are the easiest marks on the paper: learn the definitions, the tables and the diagrams in this pack.

\subsection*{How often this topic is examined}
\begin{tabularx}{\linewidth}{@{}Xll@{}}
\toprule
\textbf{Question asked} & \textbf{Paper(s)} & \textbf{Marks}\\
\midrule
Network topologies (describe 3 or 4) & 2025 Q2c, ISA2024 Q2c & 9--12\\
LAN / MAN / WAN, categories of networks & ISA2024 Q3b, 2018 Q1c & 6--9\\
Client-server vs peer-to-peer (incl.\ benefits, HIT example) & 2025 Q1b, SE2024 Q4b, Q5d & 4--8\\
Components of a data communication system (diagram) & 2018 Q1a, SE2024 Q2f & 5\\
What is a protocol (with examples) & 2025 Q2a & 5\\
Transmission modes; half duplex; point-to-point & 2023 Q4e, SE2024 Q1a.i, Q1a.ii & 1--6\\
Network design: 5 floors $\times$ 20 PCs & SE2024 Q3e & 5\\
Configurations by relationship among nodes & SE2024 Q3d & 4\\
Hardware, software \& protocols of a network & SE2024 Q4d & 4\\
Internet vs World Wide Web & SE2024 Q4a & 4\\
Ring vs bus limitations; logical vs physical topology & SE2024 Q2e, Q3c.i & 2--3\\
Network administrator / engineer competencies & SE2024 Q4e.i, Q4e.ii & 1--2\\
\bottomrule
\end{tabularx}

\begin{examtip}
The ISA paper says ``illustrate your answer, where appropriate, with large clearly labelled diagrams''.
\textbf{Draw every topology you describe} and every comparison you can. A small sketch is quick and earns marks.
\end{examtip}

%=====================================================================
\section{Data communication: components and characteristics}

\textbf{Data communication} is the exchange of data between two devices over a transmission medium.
It needs \textbf{hardware} (sender, receiver, intermediate devices) and \textbf{software} (the rules, i.e.\ the \textbf{protocol}).

\begin{center}
\begin{tikzpicture}
  \node[node,fill=cream,minimum width=30mm] (ps) at (0,1.9) {\textbf{5. Protocol}\\{\scriptsize Rule 1, Rule 2 \dots\ Rule $n$}};
  \node[node,fill=cream,minimum width=30mm] (pr) at (10,1.9) {\textbf{5. Protocol}\\{\scriptsize Rule 1, Rule 2 \dots\ Rule $n$}};
  \draw[<->,dashed,midgrey,thick] (ps) -- node[above,lbl]{both sides follow the \textbf{same} rules} (pr);
  \node[node,fill=peach,minimum width=30mm,minimum height=10mm] (s) at (0,0) {\textbf{2. Sender}\\{\scriptsize PC, phone, camera}};
  \node[node,fill=peach,minimum width=30mm,minimum height=10mm] (r) at (10,0) {\textbf{3. Receiver}\\{\scriptsize server, printer, TV}};
  \draw[->,line width=1.8pt,brand] (s) -- (r);
  \node[node,fill=white,draw=brand,minimum width=26mm] (m) at (5,0) {\textbf{1. Message}\\{\scriptsize text, numbers, images,}\\{\scriptsize audio, video}};
  \node[lbl,below=1mm of m] {\textbf{4. Transmission medium} (the path): twisted pair, coaxial, fibre optic, radio waves};
  \draw[midgrey] (ps) -- (s); \draw[midgrey] (pr) -- (r);
\end{tikzpicture}\\
\figcap{Figure: the five components of a data communication system. M-S-R-T-P.}
\end{center}

\begin{keydef}{The four characteristics of an effective data communication system}
\textbf{Delivery}: data reaches the correct destination and user. \quad
\textbf{Accuracy}: delivered without errors.\\
\textbf{Timeliness}: delivered on time (audio/video need real-time delivery). \quad
\textbf{Jitter}: the variation in packet arrival time; uneven jitter spoils real-time audio/video.
\end{keydef}
\sidenote{Lecturer's 4 characteristics: \textbf{D-A-T-J}, Delivery, Accuracy, Timeliness, Jitter. Do not confuse them with the 5 components.}

\begin{pastq}{2018 Q1a \textnormal{/} SE2024 Q2f}{5}With the aid of a fully labelled diagram, explain all the components involved in a data communication system. / With the aid of diagram, describe the complete computer communication systems.\end{pastq}
\begin{ans}
Draw the figure above, then one line per component:
\begin{enumerate}
\item \textbf{Message}: the information to be communicated: text, numbers, images, audio or video.
\item \textbf{Sender}: the device that sends the message, e.g.\ a computer, phone or video camera.
\item \textbf{Receiver}: the device the message is meant for, e.g.\ a computer, printer or TV.
\item \textbf{Transmission medium}: the physical path from sender to receiver; wired (twisted pair, coaxial, fibre optic) or wireless (radio, microwave, infrared).
\item \textbf{Protocol}: an agreed set of rules that governs the communication. Without it two connected devices cannot understand each other, like two people speaking different languages.
\end{enumerate}
\end{ans}
\sidenote{``Complete computer communication system'' (SE2024) = the same 5-component diagram. You may add a modem/NIC or switch between sender and medium to show the hardware path.}

\begin{pastq}{SE2024 Q4d}{4}A basic computer network must consist of at least hardware, software and protocols. Give examples of each of these networking requirements and explain their implementation.\end{pastq}
\begin{ans}
\small
\begin{tabularx}{\linewidth}{@{}>{\bfseries}lXX@{}}
\toprule
Requirement & Examples & Implementation / role\\\midrule
Hardware & PCs, servers, NICs, UTP and fibre cables, switches, routers, wireless access points & Physically connect devices and carry signals; a switch joins devices in a LAN, a router joins networks\\
Software & Network OS (Windows Server, Linux), NIC drivers, browsers, email clients, firewall software & Manages users, shared resources and the applications that use the network\\
Protocols & TCP/IP, Ethernet, HTTP, DNS, DHCP & Rules so devices understand each other: addressing, delivery, error control\\
\bottomrule
\end{tabularx}
\end{ans}

%=====================================================================
\section{Data flow: transmission modes and line configuration}

The \textbf{transmission mode} is the direction of signal flow between two connected devices.

\begin{center}
\begin{tikzpicture}
  \foreach \y/\t in {0/Simplex,-1.5/Half duplex,-3.0/Full duplex}{
    \node[node,fill=peach,minimum width=12mm] (a) at (0,\y) {A};
    \node[node,fill=peach,minimum width=12mm] (b) at (4.2,\y) {B};
    \node[tcap,anchor=east] at (-0.9,\y) {\t};
  }
  % simplex
  \draw[msg,line width=1.4pt] (0.7,0) -- (3.5,0);
  \node[lbl,anchor=west,text width=58mm,align=left] at (5.1,0) {One direction only; B can never reply.\\ \textit{e.g.\ keyboard $\to$ PC, CPU $\to$ monitor, TV/radio broadcast}};
  % half duplex
  \draw[msg,line width=1.4pt] (0.7,-1.35) -- node[above,lbl,pos=0.5]{time 1} (3.5,-1.35);
  \draw[msg,line width=1.4pt,dashed,draw=medc] (3.5,-1.65) -- node[below,lbl,pos=0.5]{time 2} (0.7,-1.65);
  \node[lbl,anchor=west,text width=58mm,align=left] at (5.1,-1.5) {Both directions, but \textbf{not at the same time}; the sender uses the whole channel.\\ \textit{e.g.\ walkie-talkie, CB radio}};
  % full duplex
  \draw[msg,line width=1.4pt] (0.7,-2.85) -- (3.5,-2.85);
  \draw[msg,line width=1.4pt] (3.5,-3.15) -- (0.7,-3.15);
  \node[lbl,anchor=west,text width=58mm,align=left] at (5.1,-3.0) {Both directions \textbf{at the same time}; capacity shared or two paths.\\ \textit{e.g.\ telephone / mobile call, video call}};
\end{tikzpicture}\\
\figcap{Figure: the three transmission modes (data flow).}
\end{center}

\begin{pastq}{2023 Q4e}{6}The transmission mode defines the direction of signal flow between two connected devices. State the three (3) transmission modes, giving examples for each mode.\end{pastq}
\begin{ans}
\begin{enumerate}
\item \textbf{Simplex}: communication is unidirectional; one device only sends, the other only receives. Examples: keyboard to computer, computer to monitor, radio/TV broadcast.
\item \textbf{Half duplex}: both stations can send and receive, but not at the same time; when one sends the other can only receive. Examples: walkie-talkie, CB radio.
\item \textbf{Full duplex}: both stations send and receive simultaneously. Examples: telephone or mobile phone call, video call.
\end{enumerate}
\end{ans}

\begin{pastq}{SE2024 Q1a.i}{1}Explain the following term: Half Duplex Transmission.\end{pastq}
\begin{ans}
Two-way transmission in which data flows in both directions but \textbf{only one direction at a time}, e.g.\ a walkie-talkie (press to talk, release to listen).
\end{ans}

\subsection{Line configuration: point-to-point vs multipoint}
\begin{center}
\begin{tikzpicture}
  \node[hostn] (a) at (0,0) {PC};
  \node[hostn] (b) at (3,0) {Printer};
  \draw[link,brand] (a) -- node[above,lbl]{dedicated link} (b);
  \node[note] at (1.5,-0.75) {\textbf{point-to-point}: whole link for 2 devices};
  \draw[link] (6,0) -- (12,0);
  \foreach \x/\n in {6.8/A,8.6/B,10.4/C,11.6/D}{\node[hostn,minimum width=7mm] (h\n) at (\x,0.7) {\n}; \draw[link] (h\n) -- (\x,0);}
  \node[note] at (9,-0.75) {\textbf{multipoint (multidrop)}: 3+ devices share one link};
\end{tikzpicture}\\
\figcap{Figure: the two types of connection (line configuration).}
\end{center}

\begin{pastq}{SE2024 Q1a.ii}{1}Explain the following term: Point to Point communication.\end{pastq}
\begin{ans}
A \textbf{dedicated link between exactly two devices}; the whole capacity of the link is reserved for them, e.g.\ a computer cabled directly to a printer, or a leased line between two routers. (Opposite: multipoint, where three or more devices share one link.)
\end{ans}
\sidenote{Synchronous vs asynchronous transmission (SE2024 Q1a.iii) is in Pack 03.}

%=====================================================================
\section{Categories of networks (by size)}

\begin{center}
\begin{tikzpicture}[x=17mm,y=7mm]
  \draw[midgrey,->] (0,0.6) -- (7.3,0.6);
  \foreach \x/\l in {0/1 m,1/10 m,2/100 m,3/1 km,4/10 km,5/100 km,6/1000 km,7/10\,000 km}{
    \draw[midgrey] (\x,0.5) -- (\x,0.7); \node[font=\tiny,above] at (\x,0.7) {\l};}
  \draw[fill=easyc,rounded corners=2pt] (0,-0.3) rectangle (1,0.2) node[midway,font=\scriptsize\bfseries]{PAN};
  \node[lbl,anchor=west] at (1.05,-0.05) {Bluetooth headset, phone hotspot};
  \draw[fill=peach,rounded corners=2pt] (1,-1.1) rectangle (3.3,-0.6) node[midway,font=\scriptsize\bfseries]{LAN};
  \node[lbl,anchor=west] at (3.35,-0.85) {HIT computer lab, office, campus};
  \draw[fill=medc!70,rounded corners=2pt] (3.7,-1.9) rectangle (4.7,-1.4) node[midway,font=\scriptsize\bfseries]{MAN};
  \node[lbl,anchor=east] at (3.65,-1.65) {city network, e.g.\ Harare council offices};
  \draw[fill=brand,rounded corners=2pt] (5,-2.7) rectangle (7,-2.2) node[midway,font=\scriptsize\bfseries,text=white]{WAN};
  \node[lbl,anchor=east] at (4.95,-2.45) {the Internet, bank branches across Zimbabwe};
\end{tikzpicture}\\
\figcap{Figure: network categories by geographic span (log scale).}
\end{center}

\begin{pastq}{ISA2024 Q3b \textnormal{/} 2018 Q1c}{9 / 6}Describe the differences between LAN, WAN, and MAN. / There are various categories of networks. Explain any three (3).\end{pastq}
\begin{ans}
\small
\begin{tabularx}{\linewidth}{@{}>{\bfseries}p{1.9cm}XXX@{}}
\toprule
Feature & LAN (Local Area Network) & MAN (Metropolitan Area Network) & WAN (Wide Area Network)\\\midrule
Coverage & Room, building or campus (up to a few km) & A city or town (about 5--50 km); can join several LANs & Country, continent, world; joins LANs and MANs\\
Ownership & Privately owned by one organisation & Often a private provider, ISP or municipality & Public carriers / telcos (TelOne, Liquid), leased lines\\
Speed & Very high (100 Mbps -- 10 Gbps) & High & Lower; depends on the link bought\\
Errors, delay & Low error rate, low delay & Moderate & Higher error rate and propagation delay\\
Cost & Cheap, easy to set up and maintain & Expensive to implement and maintain & Very expensive; cost varies with configuration\\
Technology & Ethernet, Wi-Fi; switches & Fibre rings, metro Ethernet, WiMAX, cable TV & Routers, leased lines, MPLS, satellite, undersea cable\\
Example & HIT computer lab & Network linking a city's offices & The Internet; a bank's branch network\\
\bottomrule
\end{tabularx}
For 2018 Q1c (any three, 6 marks): define each with size + ownership + one example (2 marks each). Others you may use: \textbf{PAN} (about 10 m, Bluetooth), \textbf{CAN} (campus), \textbf{WLAN} (wireless LAN), \textbf{SAN} (storage).
\end{ans}
\sidenote{``Describe the differences'' = a \textbf{table}. 9 marks: aim for 7 rows $\times$ 3 networks.}

\begin{pastq}{SE2024 Q4a}{4}Distinguish between the internet and the World Wide Web.\end{pastq}
\begin{ans}
\begin{tabularx}{\linewidth}{@{}XX@{}}
\toprule
\textbf{Internet} & \textbf{World Wide Web (WWW)}\\\midrule
The global network of interconnected networks: the hardware and infrastructure & A service (application) that runs \textbf{on} the Internet\\
Uses TCP/IP to move packets & Uses HTTP/HTTPS, HTML and URLs\\
Carries many services: email, FTP, VoIP, gaming and the Web & A collection of linked web pages (hypertext documents)\\
Began 1969 as ARPANET & Invented by Tim Berners-Lee, 1989--1991\\
\bottomrule
\end{tabularx}
\end{ans}

%=====================================================================
\section{Network topologies}

A \textbf{topology} is the layout pattern of the interconnections of nodes and links.
\textbf{Physical topology} = the actual layout of cables and devices (how it looks).
\textbf{Logical topology} = how data actually flows between devices (how it works).

\begin{center}
\begin{minipage}[b]{0.32\linewidth}\centering
\begin{tikzpicture}
  \draw[link,line width=1.4pt] (0,0) -- (3.6,0);
  \node[term] at (-0.05,0) {}; \node[term] at (3.65,0) {};
  \foreach \x/\s in {0.5/1,1.3/-1,2.1/1,2.9/-1}{\node[pc] (n) at (\x,\s*0.75) {PC}; \draw[link] (n) -- (\x,0);}
  \node[lbl] at (1.8,1.35) {};
\end{tikzpicture}\\[2pt]
{\small\bfseries\color{brand}Bus}\\ \figcap{one backbone, terminator at each end}
\end{minipage}\hfill
\begin{minipage}[b]{0.32\linewidth}\centering
\begin{tikzpicture}
  \node[switchn,minimum width=9mm] (c) at (0,0) {SW};
  \foreach \a in {90,162,234,306,18}{\node[pc] (p\a) at (\a:1.15) {PC}; \draw[link] (c) -- (p\a);}
\end{tikzpicture}\\[2pt]
{\small\bfseries\color{brand}Star}\\ \figcap{every node to a central hub/switch}
\end{minipage}\hfill
\begin{minipage}[b]{0.32\linewidth}\centering
\begin{tikzpicture}
  \draw[link] (0,0) circle (1.1);
  \foreach \a/\l in {90/A,30/B,-30/C,-90/D,-150/E,150/F}{\node[pc,fill=white] at (\a:1.1) {\l};}
  \draw[msg] (110:0.6) arc (110:-150:0.6);
  \node[font=\tiny] at (0,0) {token};
\end{tikzpicture}\\[2pt]
{\small\bfseries\color{brand}Ring}\\ \figcap{closed loop, data one way, token}
\end{minipage}

\bigskip
\begin{minipage}[b]{0.32\linewidth}\centering
\begin{tikzpicture}
  \foreach \i in {1,...,5}{\coordinate (m\i) at (90+72*\i-72:1.1);}
  \foreach \i in {1,...,5}{\foreach \j in {1,...,5}{\ifnum\i<\j \draw[link,thin] (m\i) -- (m\j);\fi}}
  \foreach \i/\l in {1/A,2/B,3/C,4/D,5/E}{\node[pc] at (m\i) {\l};}
\end{tikzpicture}\\[2pt]
{\small\bfseries\color{brand}Mesh (full)}\\ \figcap{$n(n-1)/2$ links: 5 nodes = 10}
\end{minipage}\hfill
\begin{minipage}[b]{0.32\linewidth}\centering
\begin{tikzpicture}
  \node[switchn,minimum width=10mm] (r) at (1.9,1.3) {root};
  \node[switchn,minimum width=8mm] (s1) at (0.95,0.3) {SW};
  \node[switchn,minimum width=8mm] (s2) at (2.85,0.3) {SW};
  \draw[link] (r) -- (s1); \draw[link] (r) -- (s2);
  \foreach \x in {0.3,0.95,1.6}{\node[pc] (q) at (\x,-0.7) {PC}; \draw[link] (s1) -- (q);}
  \foreach \x in {2.2,2.85,3.5}{\node[pc] (q) at (\x,-0.7) {PC}; \draw[link] (s2) -- (q);}
\end{tikzpicture}\\[2pt]
{\small\bfseries\color{brand}Tree (hierarchical)}\\ \figcap{stars joined to a root/backbone}
\end{minipage}\hfill
\begin{minipage}[b]{0.32\linewidth}\centering
\begin{tikzpicture}
  \draw[link,line width=1.4pt] (0,0) -- (3.6,0);
  \node[term] at (-0.05,0) {}; \node[term] at (3.65,0) {};
  \foreach \x in {0.9,2.7}{
    \node[switchn,minimum width=8mm] (s) at (\x,0.75) {SW}; \draw[link] (s) -- (\x,0);
    \foreach \d in {-0.6,0,0.6}{\node[pc] (q) at (\x+\d,1.7) {PC}; \draw[link] (s) -- (q);}}
  \node[pc] (z) at (1.8,-0.7) {PC}; \draw[link] (z) -- (1.8,0);
\end{tikzpicture}\\[2pt]
{\small\bfseries\color{brand}Hybrid (star-bus)}\\ \figcap{two or more topologies combined}
\end{minipage}\\[3pt]
\figcap{Figure: the six physical topologies.}
\end{center}

\begin{pastq}{2025 Q2c \textnormal{/} ISA2024 Q2c}{12 / 9}Describe any four (4) network topologies that can be used to connect devices. / Identify and explain any three (3) types of network topologies.\end{pastq}
\begin{ans}
\small
For each topology: sketch + how it works + one advantage + one disadvantage (3 marks each).
\smallskip

\begin{tabularx}{\linewidth}{@{}>{\bfseries}p{1.4cm}XXX@{}}
\toprule
Topology & How it works & Advantages & Disadvantages\\\midrule
Bus & All devices attach by drop cables to one backbone cable, terminated at both ends. A signal travels both ways and every node checks the address & Cheap; least cable; easy to add a node & A break in the backbone stops the whole network; collisions; hard to troubleshoot; slows as nodes are added\\
Star & Each device has its own point-to-point cable to a central hub or switch; all traffic passes through it & Easy to install and wire; a cable fault affects one node only; easy fault isolation; add/remove without disruption & Hub/switch is a single point of failure; more cable and cost than bus\\
Ring & Each device connects to its two neighbours, forming a closed loop; data travels one way and each node repeats it (Token Ring, FDDI) & No collisions (token passing); equal access; performs well under heavy load & One break or one failed node can stop the ring; adding a node disrupts it; delay grows with nodes\\
Mesh & Every device has a dedicated point-to-point link to every other: $n(n-1)/2$ links, $n-1$ ports each & Very reliable, many paths; no shared-link traffic problems; private and secure & Very expensive; lots of cable and I/O ports; hard to install\\
Tree & A hierarchy: a root node links to second-level nodes, which link to the next level (stars joined to a backbone) & Scalable; easy to manage in sections; point-to-point wiring per segment & If the root/backbone fails, segments are cut off; more difficult to configure and wire\\
Hybrid & Two or more topologies combined, e.g.\ star-bus, star-ring & Flexible; uses the strengths of each; joins different LANs & Complex to design and manage; expensive\\
\bottomrule
\end{tabularx}
\end{ans}

\begin{mnemonic}
\textbf{Mesh maths}: links $= \dfrac{n(n-1)}{2}$, ports per device $= n-1$. 5 nodes $\to$ 10 links, 4 ports each. \quad
Choosing a topology: \textbf{cost, cable length, future growth, cable type, fault tolerance}.
\end{mnemonic}

\begin{pastq}{SE2024 Q2e}{3}Explain the limitations of a ring network topology compared to a bus network topology.\end{pastq}
\begin{ans}
\begin{enumerate}
\item \textbf{Single break / failed node stops the ring}: data must pass through every node, so one failure breaks the loop (unless a dual ring is used). On a bus a failed node does not affect the others.
\item \textbf{Adding or removing a device disrupts the network}: the ring must be opened. On a bus you simply tap into the backbone.
\item \textbf{More delay and complexity}: every node must receive and regenerate the data, so delay grows with the number of nodes, and token management makes it costlier than a bus.
\end{enumerate}
\end{ans}

\begin{pastq}{SE2024 Q3c.i}{2}Distinguish between the following terms: Logical verses Physical topology.\end{pastq}
\begin{ans}
\textbf{Physical topology} is the actual layout of the cables, devices and their locations (how the network looks).
\textbf{Logical topology} is the way data flows between devices (how the network works).
Example: Ethernet through a hub is a \textbf{physical star} but a \textbf{logical bus}; Token Ring wired to a MAU is a physical star but a logical ring.
\end{ans}
\sidenote{The hub example scores the second mark: the same network can have different physical and logical topologies.}

%=====================================================================
\section{Network structure: backbone and local-access design}

Large networks are designed in two levels:
\begin{itemize}
\item \textbf{Backbone network}: the high-capacity core that joins the sites (or floors, or buildings). Few nodes (routers, core switches), fast links (fibre, leased lines), usually a \textbf{mesh or partial mesh} so there are alternative paths.
\item \textbf{Local-access network}: connects the end users' terminals to the nearest backbone node, usually a \textbf{star or tree} through switches, concentrators or multiplexers. Many cheap, lower-speed links.
\end{itemize}

\begin{center}
\begin{tikzpicture}
  \node[routern] (r1) at (4,0.9) {R1};
  \node[routern] (r2) at (7,0.9) {R2};
  \node[routern] (r3) at (4,-0.9) {R3};
  \node[routern] (r4) at (7,-0.9) {R4};
  \begin{scope}[on background layer]
    \node[draw=brand,dashed,fill=brand!8,rounded corners=8pt,fit=(r1)(r2)(r3)(r4),inner sep=4mm] (bb) {};
  \end{scope}
  \foreach \a/\b in {r1/r2,r1/r3,r2/r4,r3/r4,r1/r4}{\draw[link,brand,line width=1.4pt] (\a) -- (\b);}
  \node[lbl,above=0mm of bb] {\textbf{Backbone}: partial mesh, fibre / leased lines};
  \foreach \r/\x/\y/\d in {r1/1.9/0.9/-1,r3/1.9/-0.9/-1,r2/9.1/0.9/1,r4/9.1/-0.9/1}{
    \node[switchn,minimum width=9mm] (s) at (\x,\y) {SW};
    \draw[link] (s) -- (\r);
    \foreach \dy in {0.4,-0.1,-0.6}{\node[pc] (h) at (\x+\d*1.15,\y+\dy+0.1) {PC}; \draw[link,thin] (s) -- (h);}}
  \node[note] at (1.4,-1.95) {local access (star)}; \node[note] at (9.6,-1.95) {local access (star)};
\end{tikzpicture}\\
\figcap{Figure: a two-level design. Users reach their local switch; sites talk over the backbone.}
\end{center}

\begin{pastq}{SE2024 Q3e}{5}Identify appropriate network topology and network connecting devices for following requirements. Draw a network design for proposed network for the following scenario. An organization having its office in a building of 5 floors. Each floor needs 20 machines.\end{pastq}
\begin{ans}
\begin{itemize}
\item \textbf{Topology}: extended star (hierarchical/tree). A star on each floor; each floor switch up-links to a core switch over the vertical backbone (riser).
\item \textbf{Devices}: a 24-port access switch per floor; a core switch; a router to the ISP; a firewall; wireless access points; Cat6 UTP to desks and fibre for the riser; servers in the ground-floor server room.
\item \textbf{Addressing}: one VLAN/subnet per floor, addresses by DHCP.
\item \textbf{Why}: easy to manage and expand; a failed floor switch only affects that floor.
\end{itemize}
\begin{working}[Sizing]
\begin{align*}
\text{machines} &= 5 \times 20 = 100\\
\text{ports per floor} &= 20 + 1\ (\text{AP}) + 1\ (\text{uplink}) = 22 \le 24 && \cmt{so a 24-port switch fits, 2 spare}\\
/27 &\Rightarrow 2^{5}-2 = 30 \text{ hosts per floor} && \cmt{20 PCs + growth}
\end{align*}
\end{working}
\begin{center}
\begin{tikzpicture}[yscale=1.05]
  \foreach \f in {1,...,5}{
    \pgfmathsetmacro{\y}{\f-5}
    \draw[midgrey,dashed] (0,\y-0.5) -- (7.2,\y-0.5);
    \node[tcap,font=\scriptsize\bfseries,anchor=west] at (0,\y) {Floor \f};
    \node[hostn,minimum width=22mm] (pc\f) at (2.6,\y) {20 PCs + Wi-Fi AP};
    \node[switchn] (sw\f) at (5.4,\y) {24-port SW};
    \draw[link] (pc\f) -- (sw\f);
    \draw[link,brand] (sw\f) -- (7.4,\y);
    \node[font=\tiny,color=midgrey,anchor=south,inner sep=1pt] at (pc\f.north) {VLAN \f0: 192.168.\f0.0/27};
  }
  \draw[link,brand,line width=1.6pt] (7.4,0) -- (7.4,-5.3);
  \node[note,anchor=west,rotate=90] at (7.75,-3.6) {fibre / Cat6 riser};
  \node[tcap,font=\scriptsize\bfseries,anchor=west] at (0,-5.7) {Ground};
  \node[hostn,minimum width=22mm] (srv) at (2.6,-5.7) {DHCP / file servers};
  \node[switchn,fill=ink] (core) at (5.4,-5.7) {CORE SW};
  \draw[link] (srv) -- (core); \draw[link,brand,line width=1.6pt] (core) -- (7.4,-5.7) -- (7.4,-5.3);
  \node[node,fill=hardc!30,minimum width=12mm,font=\scriptsize] (fw) at (8.9,-5.7) {Firewall};
  \node[routern] (rt) at (10.4,-5.7) {Router};
  \node[draw=ink,fill=cream,ellipse,font=\scriptsize] (net) at (11.95,-5.7) {Internet};
  \draw[link] (7.4,-5.7) -- (fw); \draw[link] (fw) -- (rt); \draw[link] (rt) -- (net);
  \node[lbl,anchor=west,text width=38mm,align=left] at (8.4,-2) {\textbf{Extended star (tree)}:\\ a star on each floor, all floor switches up-linked to one core switch.};
\end{tikzpicture}\\
\figcap{Figure: network design for a 5-floor building, 20 machines per floor.}
\end{center}
\end{ans}

%=====================================================================
\section{Relationships among nodes: client-server and peer-to-peer}

\begin{center}
\begin{minipage}[b]{0.48\linewidth}\centering
\begin{tikzpicture}
  \node[node,fill=brand,text=white,minimum width=16mm] (s) at (0,0) {Server};
  \foreach \a/\i in {200/1,250/2,290/3,340/4}{
    \node[hostn] (c\i) at (\a:2.1) {Client};
    \draw[msg] (c\i.north) to[bend left=12] (s);
    \draw[->,thick,draw=ink] (s) to[bend left=12] (c\i.north);}
  \node[note] at (0,0.6) {stores data, gives service};
  \node[note] at (0,-2.75) {orange = request, black = response};
\end{tikzpicture}\\
{\small\bfseries\color{brand}Client-server}
\end{minipage}\hfill
\begin{minipage}[b]{0.48\linewidth}\centering
\begin{tikzpicture}
  \foreach \p/\x/\y in {1/0/0,2/2.6/0,3/0/-1.9,4/2.6/-1.9}{\node[node,fill=peach,minimum width=12mm] (p\p) at (\x,\y) {Peer \p};}
  \foreach \a/\b in {p1/p2,p1/p3,p1/p4,p2/p3,p2/p4,p3/p4}{\draw[<->,thick,draw=brand] (\a) -- (\b);}
  \node[note] at (1.3,-2.75) {each peer is both client and server};
\end{tikzpicture}\\
{\small\bfseries\color{brand}Peer-to-peer (P2P)}
\end{minipage}\\[3pt]
\figcap{Figure: the two relationship models.}
\end{center}

\begin{pastq}{2025 Q1b}{8}With the aid of proper examples, differentiate between the Client-Server model and the Peer-to-Peer model for network applications.\end{pastq}
\begin{ans}
In the \textbf{client-server} model dedicated servers provide services that clients request (the web: a browser asks a web server for a page).\par
In the \textbf{peer-to-peer} model every node is equal and both requests and provides resources (BitTorrent: each peer downloads pieces from, and uploads pieces to, other peers).\par
\small
\begin{tabularx}{\linewidth}{@{}>{\bfseries}p{1.9cm}XX@{}}
\toprule
Point & Client-server & Peer-to-peer (P2P)\\\midrule
Structure & Dedicated, powerful server(s) provide services; clients request them & All nodes are equal peers; each acts as client and server\\
Control \& security & Centralised: user accounts, access control on the server & Each user manages own machine; weaker security\\
Data storage & Stored centrally, easy backup & Scattered across peers; backup difficult\\
Scalability & Scales by adding/upgrading servers; server can become a bottleneck & Small office: slows as peers grow; file-sharing P2P scales because each peer adds capacity\\
Cost & Expensive: server hardware, network OS, administrators & Cheap: no dedicated server\\
Failure & Server is a single point of failure & No single point of failure\\
Examples & Web (browser $\leftrightarrow$ web server), email (Outlook $\leftrightarrow$ mail server), HIT e-learning portal, online banking & BitTorrent, Windows workgroup file sharing, blockchain/Bitcoin, early Skype\\
\bottomrule
\end{tabularx}
\end{ans}
\sidenote{The question says ``with the aid of proper examples''. Name real applications on both sides or you cap your mark.}

\begin{pastq}{SE2024 Q4b}{4}State the four benefits of client-server networks compared to peer to peer networks.\end{pastq}
\begin{ans}
\begin{enumerate}
\item \textbf{Centralised security}: one place for user authentication and access control.
\item \textbf{Centralised backup and data integrity}: all data on the server, backed up in one go.
\item \textbf{Easier management}: software updates, resources and users managed from the server.
\item \textbf{Better scalability and performance}: dedicated servers handle many clients; add servers as demand grows.
\end{enumerate}
\end{ans}

\begin{pastq}{SE2024 Q3d}{4}Describe in detail four possible configurations of computer networks as categorized by relationships among nodes.\end{pastq}
\begin{ans}
\begin{enumerate}
\item \textbf{Peer-to-peer}: all computers are equal; each shares its own resources and acts as client and server. No central server.
\item \textbf{Client-server}: dedicated servers (file, mail, web, database) provide services to clients that request them.
\item \textbf{Point-to-point}: a dedicated link joins exactly two nodes.
\item \textbf{Multipoint (broadcast)}: many nodes share one link or channel; what one sends is received by all (bus, Wi-Fi).
\end{enumerate}
\end{ans}
\sidenote{The wording is ambiguous. Giving P2P, client-server, point-to-point and multipoint covers both readings of ``relationship''.}

\begin{pastq}{SE2024 Q5d}{7}With reference to the client-server paradigm, explain how Harare Institute of Technology students are able to access online learning resources remotely from home.\end{pastq}
\begin{ans}
\begin{enumerate}
\item The student's laptop or phone runs a browser: this is the \textbf{client}. It joins the Internet through the home ISP (Wi-Fi or mobile data).
\item The student enters the portal URL; the client asks a \textbf{DNS server} to resolve the name to the server's IP address.
\item The client opens a \textbf{TCP connection} (3-way handshake) to the web server on port 443; \textbf{TLS} encrypts the session.
\item The client sends an \textbf{HTTP request}; the student logs in and the server \textbf{authenticates} the credentials against its user database.
\item The \textbf{server} processes the request, fetches notes, videos and assignments from central storage and sends HTTP \textbf{responses}.
\item The browser displays the resources; uploads (assignments) go the other way with HTTP POST.
\item One server serves \textbf{many clients at once}; resources are central, so every student sees the same updated material; firewalls/VPN protect the campus side.
\end{enumerate}
\begin{center}
\begin{tikzpicture}[yscale=0.6]
  \node[node,fill=peach] (c) at (0,0) {Student laptop\\{\scriptsize browser = client}};
  \node[node] (d) at (4.2,0) {DNS server};
  \node[node,fill=brand,text=white] (s) at (9.8,0) {HIT e-learning\\{\scriptsize web server + DB}};
  \foreach \x in {0,4.2,9.8}\draw[midgrey,dashed] (\x,-0.9) -- (\x,-8);
  \draw[msg] (0,-1.6) -- node[above,lbl]{1. where is the portal?} (4.2,-1.6);
  \draw[msg] (4.2,-2.6) -- node[above,lbl]{2. its IP address} (0,-2.6);
  \draw[msg,<->] (0,-3.8) -- node[above,lbl,pos=0.7]{3. TCP 3-way handshake + TLS on port 443} (9.8,-3.8);
  \draw[msg] (0,-5.0) -- node[above,lbl,pos=0.7]{4. HTTPS request: login, GET notes} (9.8,-5.0);
  \draw[->,thick,draw=ink] (9.8,-6.2) -- node[above,lbl,pos=0.3]{5. response: page, PDF, video} (0,-6.2);
  \draw[msg] (0,-7.4) -- node[above,lbl,pos=0.7]{6. POST: upload assignment} (9.8,-7.4);
\end{tikzpicture}\\
\figcap{Figure: a remote student session in the client-server model (ISP and Internet in between).}
\end{center}
\end{ans}

%=====================================================================
\section{Protocols and standards}

A \textbf{protocol} is a set of rules that governs data communication: \textbf{what} is communicated, \textbf{how} and \textbf{when}.
\begin{keydef}{Three key elements of a protocol}
\textbf{Syntax}: the structure or format of the data (e.g.\ first 8 bits = source address).\\
\textbf{Semantics}: the meaning of each section of bits and the action to take.\\
\textbf{Timing}: when data should be sent and how fast, so the receiver is not overwhelmed.
\end{keydef}

\textbf{Standards} ensure interconnectivity and interoperability, so products from different vendors work together.
\textbf{De facto} = ``by fact / convention'', adopted by widespread use without approval (e.g.\ TCP/IP grew this way).
\textbf{De jure} = ``by law / regulation'', approved by an officially recognised body.
Bodies: \textbf{ISO}, \textbf{ITU-T}, \textbf{ANSI}, \textbf{IEEE} (802.3 Ethernet, 802.11 Wi-Fi), \textbf{EIA}, \textbf{IETF} (RFCs).

\begin{pastq}{2025 Q2a}{5}Protocol is usually associated with the software components of the network, explain with the aid of examples what a protocol is?\end{pastq}
\begin{ans}
A \textbf{protocol} is an agreed set of rules that governs data communication between devices: it defines what is communicated, how and when. Without a common protocol, two connected devices are like two people speaking different languages.
Key elements: \textbf{syntax} (format of the data), \textbf{semantics} (meaning of the bits and the action to take), \textbf{timing} (when and how fast to send).
Examples:
\begin{itemize}
\item \textbf{HTTP/HTTPS}: request and deliver web pages.
\item \textbf{SMTP} (send), \textbf{POP3/IMAP} (receive): email.
\item \textbf{TCP / UDP}: reliable or fast transport; \textbf{IP}: addressing and routing.
\item \textbf{DNS}: name to IP address; \textbf{DHCP}: automatic IP configuration; \textbf{FTP}: file transfer.
\end{itemize}
\end{ans}

\begin{mnemonic}
Protocol elements: \textbf{S-S-T}, Syntax (format), Semantics (meaning), Timing (when / how fast).
\end{mnemonic}

%=====================================================================
\section{Network professionals}

\begin{pastq}{SE2024 Q4e.i \textnormal{/} Q4e.ii}{1 + 2}Write brief notes about the competencies expected of the following professionals: (i) A network administrator. (ii) A network Engineer.\end{pastq}
\begin{ans}
\textbf{(i) Network administrator} (day-to-day running): configures and monitors the network, manages user accounts and permissions, backups, patching and troubleshooting, keeps the network available.

\textbf{(ii) Network engineer} (design and build): plans and designs networks (topology, addressing, routing); selects, installs and configures routers, switches and firewalls; optimises performance and security; usually holds certifications such as CCNA/CCNP.
\end{ans}
\sidenote{Admin = \textbf{runs} it; engineer = \textbf{designs and builds} it. The hacker part (Q4e.iii) is in Pack 08.}

\begin{keydef}{Goals and applications of networks (outline week 1)}
\textbf{Goals}: resource sharing (printers, files, Internet), communication (email, video), reliability (copies, alternative paths), cost reduction, centralised management, scalability, remote access.
\textbf{Applications}: e-learning, e-banking, e-commerce, cloud services, IoT, telemedicine.
\textbf{Disadvantages}: expensive to install, needs admin time, server failure can halt work, cables break, security costs.
\end{keydef}

%=====================================================================
\section{Practice questions}
\setcounter{pq}{0}

\begin{practice}{easy}{4}State the four characteristics that determine the effectiveness of a data communication system.
\sol \textbf{Delivery} (to the correct destination), \textbf{Accuracy} (no errors), \textbf{Timeliness} (on time; real-time for audio/video), \textbf{Jitter} (variation in packet arrival time must be small).
\end{practice}

\begin{practice}{easy}{3}A company wants a full mesh between its 6 branch routers. How many links are needed and how many ports does each router need?
\sol
\begin{working}
\begin{align*}
\text{links} &= \frac{n(n-1)}{2} = \frac{6 \times 5}{2} = \mathbf{15}\\
\text{ports} &= n - 1 = \mathbf{5} \text{ per router} && \cmt{one link to each other router}
\end{align*}
\end{working}
\end{practice}

\begin{practice}{easy}{3}Classify each as simplex, half duplex or full duplex: (a) a police two-way radio (b) a WhatsApp video call (c) a car radio receiving FM.
\sol (a) Half duplex \ (b) Full duplex \ (c) Simplex.
\end{practice}

\begin{practice}{medium}{4}Distinguish between de facto and de jure standards and name two standards organisations.
\sol \textbf{De facto} (``by fact''): adopted through widespread use and convention without approval by an official body, e.g.\ TCP/IP in its early days. \textbf{De jure} (``by law''): legislated and approved by an officially recognised body, e.g.\ IEEE 802.3. Organisations: ISO, ITU-T, ANSI, IEEE, EIA, IETF.
\end{practice}

\begin{practice}{medium}{5}Explain the three elements of a protocol, using an email message as the example.
\sol \textbf{Syntax}: the format: an email has header fields (From, To, Subject) followed by a body, in a fixed order. \textbf{Semantics}: what the fields mean and the action: ``To'' tells the mail server where to deliver; a 550 reply means the mailbox does not exist, so the sender gives up. \textbf{Timing}: when and how fast: the client waits for the server's ``250 OK'' before sending the next command, so it does not overwhelm the server.
\end{practice}

\begin{practice}{medium}{6}Compare star and mesh topologies under the headings reliability, cost, installation and typical use.
\sol
\begin{tabularx}{\linewidth}{@{}>{\bfseries}lXX@{}}
\toprule
 & Star & Mesh\\\midrule
Reliability & One cable fault affects one node, but the central switch is a single point of failure & Very high: many alternative paths, no single point of failure\\
Cost & Moderate: one cable per node + a switch & Very high: $n(n-1)/2$ links and $n-1$ ports per node\\
Installation & Easy to install, add and remove nodes & Difficult; adding a node means links to all others\\
Typical use & Office and lab LANs & WAN backbones between routers\\
\bottomrule
\end{tabularx}
\end{practice}

\begin{practice}{medium}{4}A school lab uses Ethernet with all PCs cabled to a hub. State its physical and logical topology and explain why they differ.
\sol Physical: \textbf{star} (each PC has its own cable to the hub). Logical: \textbf{bus}, because a hub repeats every frame out of all ports, so every PC receives every frame and they share one collision domain, exactly as on a shared bus cable. Physical = how it is wired; logical = how the data flows.
\end{practice}

\begin{practice}{hard}{12}A retail company has a head office in Harare and three branches (Bulawayo, Mutare, Gweru), each with about 30 staff. With the aid of a diagram, propose a design using a backbone network and local-access networks. Justify the topology and devices chosen at each level.
\sol
\begin{center}
\begin{tikzpicture}
  \node[routern] (h) at (0,0) {HRE};
  \node[routern] (b) at (-3,-1.6) {BYO};
  \node[routern] (m) at (3,-1.6) {MUT};
  \node[routern] (g) at (0,-3.0) {GWE};
  \foreach \a/\b in {h/b,h/m,h/g,b/g,m/g}{\draw[link,brand,line width=1.4pt] (\a) -- (\b);}
  \node[lbl] at (1.2,-1.6) {WAN backbone\\(partial mesh)};
  \foreach \r/\x/\y in {h/0/1.25,b/-5.2/-1.6,m/5.2/-1.6,g/0/-4.25}{
    \node[switchn,minimum width=12mm] (s) at (\x,\y) {LAN SW}; \draw[link] (s) -- (\r);}
  \node[note] at (-5.2,-2.15) {star, 30 PCs}; \node[note] at (5.2,-2.15) {star, 30 PCs};
  \node[note,anchor=west] at (0.8,1.25) {HQ: servers, firewall, Internet};
\end{tikzpicture}
\end{center}
\begin{itemize}
\item \textbf{Backbone}: one router per site joined by leased lines / MPLS / fibre from a carrier, in a \textbf{partial mesh} so each site has at least two paths; if one link fails traffic reroutes. Full mesh (6 links) would cost more for little gain.
\item \textbf{Local access}: in each site a \textbf{star} of PCs on a 48-port switch (or two 24-port switches, extended star), Cat6 UTP, a wireless AP; the switch up-links to the site router.
\item \textbf{HQ}: servers (ERP, email, file, DHCP/DNS), firewall and the single Internet link; branches reach them over the backbone.
\item \textbf{Addressing}: one subnet per site, e.g.\ 10.1.0.0/24, 10.2.0.0/24 \dots
\item \textbf{Justification}: star is cheap, easy to manage, a fault affects one PC; the mesh backbone gives reliability where an outage would stop all trading; the two-level design scales by adding a site router and one more star.
\end{itemize}
\end{practice}

\begin{practice}{hard}{10}Discuss the factors a network designer should consider when choosing a topology for a new hospital network, and recommend a topology.
\sol
\begin{enumerate}
\item \textbf{Reliability / fault tolerance}: patient systems must not stop; avoid bus and ring, where one break stops all.
\item \textbf{Cost}: budget for cable, switches and ports (mesh is costly).
\item \textbf{Cable length and building layout}: several wards and floors need a hierarchy.
\item \textbf{Future growth}: new wards must be added without disrupting others; star/tree expand easily.
\item \textbf{Ease of troubleshooting}: faults must be isolated quickly.
\item \textbf{Performance}: imaging (X-ray) files need high bandwidth, no shared collision domain.
\item \textbf{Security}: separate clinical, admin and guest traffic (VLANs).
\end{enumerate}
\textbf{Recommendation}: a \textbf{hybrid}: an extended star (tree) in each building for local access, with redundant core switches joined in a \textbf{partial mesh} backbone, so no single device failure stops the hospital. Draw the core mesh with stars hanging off it.
\end{practice}

\begin{practice}{hard}{10}``Peer-to-peer networks are obsolete; every organisation should use client-server.'' Discuss this statement, with examples.
\sol
\textbf{For client-server}: centralised security and authentication; central backup and data integrity; easier management of users and updates; dedicated servers give consistent performance; essential for banks, universities (HIT e-learning), email and web services.\\
\textbf{Against (P2P still useful)}: cheap, no server or network OS to buy; quick for a small office of 2--10 PCs sharing a printer (Windows workgroup); no single point of failure; large-scale P2P (BitTorrent, blockchain) scales because every peer adds capacity and bandwidth; reduces load on central servers (e.g.\ software updates delivered peer to peer).\\
\textbf{Conclusion}: the statement is too strong. Client-server is best where security, central data and many users matter; P2P remains the right choice for very small networks and for distributed content sharing. Many systems are hybrid (a central index server + peer transfers).
\end{practice}

\begin{practice}{hard}{12}With the aid of diagrams, explain the bus, ring and star topologies, and show what happens in each when one cable segment fails.
\sol
Draw each topology (Section 4) and mark a cross on one link.
\begin{itemize}
\item \textbf{Bus}: all nodes on one terminated backbone; a signal goes both ways and every node checks the address. \textbf{Backbone break}: the cable is no longer terminated, signals reflect and the \textbf{whole segment goes down}. A drop-cable break isolates only that node.
\item \textbf{Ring}: each node connects to two neighbours; data travels one way and each node repeats it (token passing). \textbf{One break}: the loop is open, the token cannot circulate and the \textbf{whole ring stops} (a dual ring such as FDDI wraps round the fault).
\item \textbf{Star}: each node has its own link to a central switch. \textbf{One cable break}: only that node is cut off; the rest work. But if the \textbf{switch} fails, everything stops.
\end{itemize}
Conclusion: star gives the best fault isolation for a LAN, which is why modern Ethernet LANs are stars.
\end{practice}

%=====================================================================
\newpage
\begin{cheatsheet}
\small
\begin{tabularx}{\linewidth}{@{}XX@{}}
\textbf{5 components} (M-S-R-T-P)
\begin{itemize}
\item Message, Sender, Receiver, Transmission medium, Protocol
\item 4 characteristics: Delivery, Accuracy, Timeliness, Jitter
\end{itemize}
\textbf{Transmission modes}
\begin{itemize}
\item Simplex: one way (keyboard, TV broadcast)
\item Half duplex: both ways, one at a time (walkie-talkie)
\item Full duplex: both ways at once (phone call)
\item Point-to-point: dedicated link, 2 devices; multipoint: 3+ share a link
\end{itemize}
\textbf{Network sizes}
\begin{itemize}
\item PAN $\sim$10 m (Bluetooth); LAN building/campus, private, fast, cheap
\item MAN city 5--50 km; WAN country/world, carriers, slower, costly
\item Internet = network of networks (TCP/IP); WWW = web pages on it (HTTP)
\end{itemize}
\textbf{Protocol} = rules: what, how, when. Syntax, Semantics, Timing. HTTP, SMTP, TCP, IP, DNS, DHCP, FTP.\par
\textbf{Standards}: de facto (by use) vs de jure (by law). ISO, ITU-T, ANSI, IEEE, IETF.
&
\textbf{Topologies}
\begin{itemize}
\item \textbf{Bus}: one backbone + terminators; cheap; break = all down
\item \textbf{Star}: central switch; easy, fault isolation; switch = single point of failure
\item \textbf{Ring}: loop, one way, token; no collisions; one break stops ring
\item \textbf{Mesh}: $n(n-1)/2$ links, $n-1$ ports; reliable, very costly
\item \textbf{Tree}: hierarchy of stars on a root; scalable; root failure cuts segments
\item \textbf{Hybrid}: mix (star-bus); flexible, complex
\item Physical = how it is wired; logical = how data flows (hub: star wired, bus logic)
\end{itemize}
\textbf{Design}: backbone = fast core, mesh/partial mesh; local access = star/tree to users. 5 floors $\times$ 20 PCs: extended star, 24-port SW per floor, core SW, router, firewall, VLAN per floor.

\textbf{Client-server}: central server, security, backup, management, scalable; costly, server = SPOF. Web, email, HIT portal.\par
\textbf{P2P}: equal peers, cheap, no SPOF; weak security, scattered data. BitTorrent, workgroup.

\textbf{Admin} runs the network; \textbf{engineer} designs and builds it.
\end{tabularx}
\end{cheatsheet}

\end{document}
